% Options for packages loaded elsewhere
\PassOptionsToPackage{unicode}{hyperref}
\PassOptionsToPackage{hyphens}{url}
%
\documentclass[
  10pt,
]{article}
\usepackage{lmodern}
\usepackage{amssymb,amsmath}
\usepackage{ifxetex,ifluatex}
\ifnum 0\ifxetex 1\fi\ifluatex 1\fi=0 % if pdftex
  \usepackage[T1]{fontenc}
  \usepackage[utf8]{inputenc}
  \usepackage{textcomp} % provide euro and other symbols
\else % if luatex or xetex
  \usepackage{unicode-math}
  \defaultfontfeatures{Scale=MatchLowercase}
  \defaultfontfeatures[\rmfamily]{Ligatures=TeX,Scale=1}
\fi
% Use upquote if available, for straight quotes in verbatim environments
\IfFileExists{upquote.sty}{\usepackage{upquote}}{}
\IfFileExists{microtype.sty}{% use microtype if available
  \usepackage[]{microtype}
  \UseMicrotypeSet[protrusion]{basicmath} % disable protrusion for tt fonts
}{}
\makeatletter
\@ifundefined{KOMAClassName}{% if non-KOMA class
  \IfFileExists{parskip.sty}{%
    \usepackage{parskip}
  }{% else
    \setlength{\parindent}{0pt}
    \setlength{\parskip}{6pt plus 2pt minus 1pt}}
}{% if KOMA class
  \KOMAoptions{parskip=half}}
\makeatother
\usepackage{xcolor}
\IfFileExists{xurl.sty}{\usepackage{xurl}}{} % add URL line breaks if available
\IfFileExists{bookmark.sty}{\usepackage{bookmark}}{\usepackage{hyperref}}
\hypersetup{
  hidelinks,
  pdfcreator={LaTeX via pandoc}}
\urlstyle{same} % disable monospaced font for URLs
\usepackage[a4paper,margin=18mm]{geometry}
\setlength{\emergencystretch}{3em} % prevent overfull lines
\providecommand{\tightlist}{%
  \setlength{\itemsep}{0pt}\setlength{\parskip}{0pt}}
\setcounter{secnumdepth}{-\maxdimen} % remove section numbering

\usepackage{esint}
\usepackage{xeCJK}
\setmainfont{DejaVu Serif}
\setmonofont{DejaVu Sans Mono}
\setCJKmainfont{Noto Serif CJK SC}
\setCJKmonofont{Noto Sans Mono CJK SC}
\setmathfont{Asana Math}
\setmathfont{Latin Modern Math}[range=\lrcorner]
\setlength{\emergencystretch}{3em}
\allowdisplaybreaks
\hypersetup{hypertexnames=false}

\author{}
\date{}

\begin{document}

\hypertarget{the-actual-torsion-bracket-and-its-integral-prolongation}{%
\section{The actual torsion bracket and its integral
prolongation}\label{the-actual-torsion-bracket-and-its-integral-prolongation}}

28 September 2026. This extends, and removes the exact-closure
restriction from, Section 3.2 of the
\href{20260928_zero_endpoint_phase_curvature.md}{phase-curvature proof}.
The actual constant-weight monopole equations are retained. The result
is an exact differential identity and its tested form, not the missing
weak zero-end gap.

\hypertarget{the-bracket-is-determined-by-the-actual-cubic-torsion}{%
\subsection{1. The bracket is determined by the actual cubic
torsion}\label{the-bracket-is-determined-by-the-actual-cubic-torsion}}

Use the Clifford matrices from I41, Section 14,

\[
 C_1=I,\quad C_2=\operatorname{diag}(i,-i),\quad
 C_3=\begin{pmatrix}0&1\\-1&0\end{pmatrix},\quad
 C_4=\begin{pmatrix}0&-i\\-i&0\end{pmatrix}.
\]

On the open set rho = \textbar Phi\textbar{} squared \textgreater{} 0,
choose a local positive-spinor frame with \(\Phi=\sqrt\rho\,e_1\). All
positive and negative spinor frames, the tangent frame and the
trivialization of L are chosen compatibly with the displayed Clifford
matrices. Further adaptations below change these compatible frames
together. Write

\[
 \nabla_j\Phi=(h_j,k_j),\qquad
 \zeta_j=\langle\Phi,\nabla_j\Phi\rangle=\sqrt\rho\,h_j.
\]

The defining cubic torsion \(\tau=t_\Phi\), a section of
\(E=S^-\otimes L\), has components

\[
             \tau=\rho(k_1+ik_2,-k_3-ik_4).                \tag{1.1}
\]

Let \(\Omega_+=e^{12}+e^{34}\), \(\Psi_2=e^{13}-e^{24}\),
\(\Psi_3=e^{14}+e^{23}\). The off-diagonal positive-spinor connection is
the complex one-form \(\eta_j=k_j/\sqrt\rho\). The Clifford action gives

\[
 \nabla\Omega_+=2(\operatorname{Im}\eta)\otimes\Psi_2
                         -2(\operatorname{Re}\eta)\otimes\Psi_3.
\]

This formula first concerns real one-forms; evaluation on complex
vectors is then extended complex linearly.

Put \(Z_1=e_1-ie_2\), \(Z_2=e_3-ie_4\). Consequently

\[
 (\nabla_XZ_1)^{0,1}=\overline\eta(X)\overline Z_2,
 \qquad (\nabla_XZ_2)^{0,1}=-\overline\eta(X)\overline Z_1.
\]

Here a conjugated one-form is evaluated complex linearly; in particular
\(\overline\eta(Z_1)\) is not \(\overline{\eta(Z_1)}\). Torsion-freeness
of the Levi-Civita connection and (1.1) give

\[
 \boxed{[Z_1,Z_2]^{0,1}
  =-\frac{\overline\tau_1}{\rho^{3/2}}\overline Z_1
    +\frac{\overline\tau_2}{\rho^{3/2}}\overline Z_2.}      \tag{1.2}
\]

This identity uses the cubic definition, not closure of Omega or the
Dirac equation. It is smooth through zeros of tau on rho \textgreater{}
0. It is a local expression of a tensor identity, so changing adapted
frames does not change its content. No assertion across rho = 0 is
needed below.

On rho,m \textgreater{} 0, where \(m=|\tau|^2\), further adapt the
negative frame so that \(\tau=(\sqrt m,0)\), and set
\(\kappa_0=\sqrt m/(2\rho^{3/2})\). Then (1.2) becomes exactly

\[
                [Z_1,Z_2]^{0,1}=-2\kappa_0\overline Z_1.   \tag{1.3}
\]

One may also see the cancellation directly in the actual Dirac system.
Its two rows are

\[
 h_1+ih_2+k_3-ik_4=0,\qquad
 -h_3-ih_4+k_1-ik_2=0.
\]

Solving together with (1.1) yields

\[
 \eta=(\kappa_0+z,-i\kappa_0+iz,-w,-iw),\quad
 z=\frac{\zeta_3+i\zeta_4}{2\rho},\quad
 w=\frac{\zeta_1+i\zeta_2}{2\rho}.
\]

The z and w contributions cancel from \(\overline\eta(Z_1)\) and
\(\overline\eta(Z_2)\). They need not vanish from other connection
components. Thus there is no separate near-closure correction to (1.3).

\hypertarget{the-curvature-forcing-is-exactly-the-trace-free-ricci-block}{%
\subsection{2. The curvature forcing is exactly the trace-free Ricci
block}\label{the-curvature-forcing-is-exactly-the-trace-free-ricci-block}}

For the actual constant-weight solution, \(F_A^+\) is a multiple of
\(i\rho\Omega_+\). All anti-self-dual forms have type (1,1), and
\(Z_1\wedge Z_2=\Psi_2-i\Psi_3\) is self-dual and perpendicular to
Omega. Hence \(F_A(Z_1,Z_2)=0\), exactly.

The connection on E has central curvature \((3/2)F_A I\). Negative
Clifford multiplication annihilates self-dual forms, so the block
decomposition of the Riemannian curvature operator gives

\[
 \boxed{F^E(Z_1,Z_2)
   =-\tfrac12c^-\bigl(B_{\operatorname{Ric}^{\circ}}^*
                          (\Psi_2-i\Psi_3)\bigr).}       \tag{2.1}
\]

Here \(B_{\operatorname{Ric}^{\circ}}^*=\pi_-\mathcal R|_{\Lambda^+}\)
and \((F_{\rm LC}^{S^-})_{ij}=-\tfrac12c^-(\mathcal R(e^i\wedge e^j))\),
as in the preceding phase-curvature proof.

In particular
\(|F^E(Z_1,Z_2)\tau|\le C|\operatorname{Ric}^{\circ}|\sqrt m\). Neither
Weyl block nor an independent determinant error occurs in this curvature
component. This does not control all components of F\^{}E.

Define \(A_i=\nabla_{Z_i}\tau\), \(B=\nabla_{\overline Z_1}\tau\), and
let \(\sum L_iZ_i=[Z_1,Z_2]^{1,0}\). Commuting actual derivatives and
using (1.3) gives

\[
 \boxed{2\kappa_0 B=F^E(Z_1,Z_2)\tau
       -\nabla_{Z_1}A_2+\nabla_{Z_2}A_1+\sum_i L_iA_i.}   \tag{2.2}
\]

This is the full actual counterpart of the earlier ideal-system
calculation. The moving-frame terms L\_i remain part of the equation.
Smallness of A\_i in L2 is not assumed: I41 (14.2) itself retains an
actual current--Hessian term in their projected derivative.

\hypertarget{testing-does-not-remove-the-symmetric-second-derivatives}{%
\subsection{3. Testing does not remove the symmetric second
derivatives}\label{testing-does-not-remove-the-symmetric-second-derivatives}}

Let \(\mathcal A=(\nabla\tau)^{1,0}\), regarded as an E-valued one-form.
Let \(\mu\) denote the algebraic almost-complex component of exterior
covariant differentiation from type (0,1) to type (2,0). In the frame
above, \((\mu(\nabla\tau)^{0,1})(Z_1,Z_2)=2\kappa_0 B\). The invariant
form of (2.2) is

\[
 \partial_\nabla\mathcal A+\mu(\nabla\tau)^{0,1}
                         =(F^E)^{2,0}\tau.               \tag{3.1}
\]

For a real smooth cutoff f with compact support in rho,m \textgreater{}
0, integration by parts is legitimate and gives

\[
\begin{aligned}
 \int f^2|\mu(\nabla\tau)^{0,1}|^2
 ={}&\operatorname{Re}\int f^2
       \langle(F^E)^{2,0}\tau,\mu(\nabla\tau)^{0,1}\rangle\\
 &-\operatorname{Re}\int\left\langle\mathcal A,
       \partial_\nabla^*\{f^2\mu(\nabla\tau)^{0,1}\}\right\rangle.
\end{aligned}                                             \tag{3.2}
\]

The adjoint contains derivatives of the cutoff, frame, mu and
\(\nabla\tau\). Its principal second derivatives have not acquired a
sign. This formula does remove an unnecessary independent closure error,
but it does not give an L2 estimate on B from an L2 estimate on A. All
norms here are tensor norms: Z\_i have squared length two, and their
dual (1,0)-forms have squared norm one half. Thus
\(|\mathcal A|^2=(|A_1|^2+|A_2|^2)/2\) and
\(|\mu(\nabla\tau)^{0,1}|^2=\kappa_0^2|B|^2\). Evaluating (3.1) on
Z\_1,Z\_2 does not preserve these norm factors.

\hypertarget{the-remaining-weyl-interaction-in-actual-derivative-variables}{%
\subsection{4. The remaining Weyl interaction in actual derivative
variables}\label{the-remaining-weyl-interaction-in-actual-derivative-variables}}

Retain the opposite-Weyl frame of the phase-curvature proof. Set
\(D=(W_{22}-W_{33})/2\), \(E_0=W_{23}\), and
\(C_0=\nabla_{\overline Z_2}\tau\). Subscripts 1,2 now denote spinor
components, not tangent indices. The free planar contribution to
\(\operatorname{Re}\langle\nabla\tau\wedge  F_{\rm LC}^{S^-}(W^-)\wedge\nabla\tau\rangle\)
is exactly

\[
 \boxed{\operatorname{Re}\left[(D-iE_0)
    \left\{\overline{B_1}(C_0)_2
       -\overline{(A_2)_1}(A_1)_2\right\}\right]d\mu.}   \tag{4.1}
\]

There is no N = 0 or torsion-Dirac = 0 hypothesis here. To verify it,
write \(\nabla\tau=(x,y)\) as two complex one-forms. The old wedge
matrix used derivatives divided by square root m; here that factor is
absorbed: its second row plus i times its third row, multiplied by m, is
the contraction of \(\overline x\wedge y\). Since

\[
 \Omega_2+i\Omega_3=(e^1-ie^2)\wedge(e^3+ie^4),
\]

evaluate directly \((\overline x\wedge y)(Z_1,\overline Z_2)\). It
equals \(\overline{B_1}(C_0)_2-\overline{(A_2)_1}(A_1)_2\), proving
(4.1). This also reproduces the earlier \(4m\{D(au+bv)+E_0(av-bu)\}\)
under that calculation's additional pointwise constraints.

Substituting (2.2) into (4.1) does not turn it into a nonnegative
curvature square. Integration by parts differentiates the opposite Weyl
block, the cutoff and C\_0. In particular it leaves symmetric barred
second derivatives of tau, not just curvature commutators. These are
precisely the quantities that a further signed integral identity would
have to absorb. No estimate for them follows from separate L2 bounds for
curvature and the first derivative of tau.

\hypertarget{a-finite-jet-check-of-the-surviving-derivative}{%
\subsubsection{4.1 A finite jet check of the surviving
derivative}\label{a-finite-jet-check-of-the-surviving-derivative}}

This check locates a genuine principal symbol; it is not a closed
solution or a theorem ruling out additional global cancellations. At a
point put \(z_1=x_1+ix_2,z_2=x_3+ix_4\) and retain the constant Clifford
matrices. For a positive spinor increment (H,G), the two Dirac rows are

\[
 D^+(H,G)=\bigl((\partial_1+i\partial_2)H+
                   (\partial_3-i\partial_4)G,
       (-\partial_3-i\partial_4)H+(\partial_1-i\partial_2)G\bigr).
\]

The linearized torsion at \(\Phi=\sqrt\rho e_1\) is
\(T(H,G)=\rho((\partial_1+i\partial_2)G,  - (\partial_3+i\partial_4)G)\).
Both polynomial seeds

\[
 (H,G)=(\overline z_1\overline z_2,|z_1|^2-|z_2|^2),
 \qquad (H,G)=(0,\overline z_2^3/6)
\]

solve this flat D-plus system exactly. Their torsions are respectively
\(\rho(2z_1,2z_2)\) and \(\rho(0,-\overline z_2^2)\); both solve the
negative Dirac system

\[
 D^-\tau=\bigl((A_2)_2-(A_1)_1,-B_2-(C_0)_1\bigr)=0.
\]

They give \((A_1)_1=(A_2)_2=4\rho\) and
\(\nabla_{\overline Z_2}^2\tau_2=-8\rho\), respectively. After
substituting (2.2) in (4.1), the adjoint retains the term

\[
 \frac{f}{2\kappa_0}\operatorname{Re}\left[(D-iE_0)
 \left\{\overline{(A_2)_1}\nabla_{\overline Z_1}
                          \nabla_{\overline Z_2}\tau_2
       -\overline{(A_1)_1}\nabla_{\overline Z_2}^2\tau_2\right\}\right],
\]

in addition to curvature, derivative-Weyl, cutoff and frame terms. Here
f denotes the scalar weight of that integration. Reversing the cubic
seed reverses this displayed pairing while leaving the lower jets
unchanged. Its fully symmetric third-spinor jet also leaves the
principal curvature-commutator constraints unchanged. Thus neither the
two Dirac contractions nor commuting antisymmetric derivatives makes
this symmetric symbol vanish pointwise. The exact integral formula
(3.2), not this isolated symbol, is the global identity to use.

\hypertarget{the-global-condition-to-which-the-calculation-must-return}{%
\subsection{5. The global condition to which the calculation must
return}\label{the-global-condition-to-which-the-calculation-must-return}}

The same-solution characteristic localization in I41, Section 7, is

\[
 16\pi^2\int d\chi\wedge T_\tau
     =2\int\vartheta^2|W^-|^2+\int\vartheta^2 k,
\] \[
 k=9(|F_A^+|^2-|F_A^-|^2)+s^2/24
                                -|\operatorname{Ric}^{\circ}|^2/2.
\]

Here \(\chi=1-\vartheta^2\) cuts off the fixed low-torsion transition
band. Along a putative actual zero-end sequence, the last integral
divided by V tends to 20/3. A sufficient new estimate is that the net
flux after subtracting the displayed opposite-Weyl payment has upper
limit strictly below \(20V/3\); an o(V) bound would suffice but is
stronger than necessary. Such a bound is not proved by (3.2) or (4.1).

The update is therefore precise: exact closure need not be added to the
actual bracket, and its (2,0) curvature forcing is only Ricci error. The
symmetric derivative interaction and its global signed absorption remain
unresolved. O10 stays open with the original actual supply, class-square
and normalization requirements.

\end{document}
